% Part: many-valued-logic
% Chapter: sequent-calculus
% Section: introduction

\documentclass[../../../include/open-logic-section]{subfiles}

\begin{document}

\olfileid{mvl}{seq}{int}

\olsection{مقدمة}

حساب التتابعيات في المنطق الكلاسيكي نسق !!{derivation} يجمع البساطة
إلى الفعالية. ولا يختص إمكان بنائه بذلك المنطق؛ فمتى عُرّف منطق متعدد
القيم بمصفوفة كانت مجموعة قيمها~$\OLId{latin-looped}{V}$ منتهية، أمكن أن يُجعل له حساب
تتابعيات. ولبيان الطريق إلى ذلك، ننظر أولًا في دلالة التتابعية الكلاسيكية
وفي صورة قواعد الاستدلال عليها.

تذكر الآن أن التتابعية
\begin{align*}
    !A_\OLMathNumeral{1}, \dots, !A_\OLId{latin-ordinary}{m} & \Sequent !B_\OLMathNumeral{1}, \dots, !B_\OLId{latin-ordinary}{n}
\intertext{يمكن تفسيرها بالـ!!{formula}}
    (!A_\OLMathNumeral{1} \land \cdots \land !A_\OLId{latin-ordinary}{m}) & \lif (!B_\OLMathNumeral{1} \lor \cdots \lor
!B_\OLId{latin-ordinary}{n})
\end{align*}
ومعنى ذلك أن !!^a{valuation}~$\pAssign{\OLId{latin-ordinary}{v}}$ \emph{يحقق} التتابعية
$\OLId{greek-variable}{Gamma} \Sequent \OLId{greek-variable}{Delta}$ متى وفقط متى حصل أحد الأمرين:
$\pValue{\OLId{latin-ordinary}{v}}(!A) = \False$ لبعض~$!A \in \OLId{greek-variable}{Gamma}$، أو
$\pValue{\OLId{latin-ordinary}{v}}(!A) = \True$ لبعض~$!A \in \OLId{greek-variable}{Delta}$. وعلى هذا تكون
التتابعيات الابتدائية~$!A \Sequent !A$ محققة على كل حال؛ إذ لا يخلو
الأمر من~$\pValue \OLId{latin-ordinary}{v}(!A) = \True$ أو~$\pValue \OLId{latin-ordinary}{v}(!A) = \False$.

فيما يأتي قاعدتا الاستدلال الخاصتان بالشرط في~$\Log{LK}$، بعد حذف
الصيغ الجانبية~$\OLId{greek-variable}{Gamma}$ و$\OLId{greek-variable}{Delta}$:

\begin{defish}
    \Axiom$ \fCenter !A$
    \Axiom$ !B \fCenter $
    \RightLabel{\LeftR{\lif}}
    \BinaryInf$ !A \lif !B \fCenter $
    \DisplayProof
    \hfill
    \Axiom$ !A \fCenter !B$
    \RightLabel{\RightR{\lif}}
    \UnaryInf$ \fCenter !A \lif !B $
    \DisplayProof
\end{defish}

فلنقرأ القاعدتين على مقتضى هذه الدلالة. تفيد قاعدة
$\LeftR{\lif}$ أن اجتماع~$\pValue \OLId{latin-ordinary}{v}(!A) = \True$ و
$\pValue \OLId{latin-ordinary}{v}(!B) = \False$ يقتضي~$\pValue \OLId{latin-ordinary}{v}(!A \lif !B) = \False$.
وتفيد قاعدة~$\RightR{\lif}$ أن حصول أحد الأمرين:
$\pValue \OLId{latin-ordinary}{v}(!A) = \False$ أو~$\pValue \OLId{latin-ordinary}{v}(!B) = \True$، يقتضي
$\pValue \OLId{latin-ordinary}{v}(!A \lif !B) = \True$. بل يصح العكس في كل منهما،
فتكون الدلالة تكافؤًا. وليس هذا شأن قاعدتي~$\LeftR{\land}$ و$\RightR{\lor}$
في صورتيهما القياسيتين، غير أن لكل منهما صورة أخرى يتحقق فيها التكافؤ:

\begin{defish}
    \Axiom$!A, !B, \OLId{greek-variable}{Gamma} \fCenter \OLId{greek-variable}{Delta}$
    \RightLabel{\LeftR{\land}}
    \UnaryInf$!A \land !B, \OLId{greek-variable}{Gamma} \fCenter \OLId{greek-variable}{Delta}$
    \DisplayProof
    \hfill
    \Axiom$ \OLId{greek-variable}{Gamma} \fCenter \OLId{greek-variable}{Delta}, !A, !B$
    \RightLabel{\RightR{\lor}}
    \UnaryInf$ \OLId{greek-variable}{Gamma} \fCenter \OLId{greek-variable}{Delta}, !A \lor !B$
    \DisplayProof
\end{defish}

فإذا نقلنا هذا المعنى إلى منطق ذي~$\OLId{latin-ordinary}{n}$ من القيم، كان حساب التتابعيات
ذا~$\OLId{latin-ordinary}{n}$ من المواضع، لا موضعين فحسب: لكل قيمة صدق موضع يخصها.
ومثال ذلك المنطق الثلاثي القيم حيث~$\OLId{latin-looped}{V} = \{\False, \Undef, \True\}$؛
فالتتابعية فيه على صورة~$\OLId{greek-variable}{Gamma} \mid \Pi \mid \OLId{greek-variable}{Delta}$. ويحققها
!!a{valuation}~$\pAssign \OLId{latin-ordinary}{v}$ متى وفقط متى صدق أحد الشروط الآتية:
$\pValue{\OLId{latin-ordinary}{v}}(!A) = \False$ لبعض~$!A \in \OLId{greek-variable}{Gamma}$، أو
$\pValue{\OLId{latin-ordinary}{v}}(!A) = \True$ لبعض~$!A \in \OLId{greek-variable}{Delta}$، أو
$\pValue{\OLId{latin-ordinary}{v}}(!A) = \Undef$ لبعض~$!A \in \Pi$.
ولذلك كانت التتابعيات الابتدائية~$!A \mid !A \mid !A$ محققة دائمًا.

\end{document}
